\subsection{Methodology}
This replication was conducted as an independent re-implementation from the paper's description, without access to the original authors' codebase. All experiments were re-written from scratch in Python using publicly available model checkpoints, libraries, and datasets. The sections reproduced are Section 3.3 (object discovery), Section 3.4 (qualitative evaluation), Appendix C (analysis of LOST performance), and Appendix D (behaviour of models trained with registers). Code was run on a single NVIDIA GPU, a single/double T4 (on Kaggle), with model loading handled via the \texttt{timm} and \texttt{open\_clip} libraries and, where necessary, HuggingFace \texttt{transformers}.

\subsubsection{Model descriptions}
Table~\ref{tab:models} summarises the model configurations used across the reproduced experiments.
Original checkpoints for both OpenCLIP+reg and DeiT-III+reg were never publicly released: the OpenCLIP proxy follows \citet{jiang2025vision}; the DeiT-III proxy uses \texttt{vit\_base\_patch14\_reg4\_dinov2.lvd142m}, the closest available timm checkpoint, and is clearly labelled as a proxy in all results.

\subsubsection{Datasets}
The datasets used are FGVC-Aircraft (\texttt{trainval}/\texttt{test}), Oxford-IIIT Pet (500 subsampled),
COCO val2017 (4{,}952 annotated images), PASCAL VOC 2007 (\texttt{trainval}, 5{,}011 images), and
PASCAL VOC 2012 (\texttt{trainval}, 11{,}540 images), all resized to the model's native resolution and
normalised with ImageNet mean/std. No augmentation is applied (full details in Table~\ref{tab:datasets}).

\subsubsection{Hyperparameters}
All hyperparameters were taken directly from the paper or, where not specified, from standard practice.
\textbf{LOST algorithm} (experiments: object discovery, lost intermediate maps) hyperparameters are reported in Table~\ref{tab:LOSTalgo};
\textbf{linear probing} (experiment information held by tokens) hyperparameters in Table~\ref{tab:linearprobe};
and \textbf{attention map averaging} (experiment: positional focus) hyperparameters in Table~\ref{tab:attmapavg}.\\
Three deviations from the paper are noted.
\textit{(1) Per-row mean-centering:} without centering, DINOv2 keys produce an all-positive Gram matrix
causing seed selection to collapse to a fixed corner patch; we subtract the per-row mean from $G$ after
bias addition. This is the primary cause of the systematic 15--20\% CorLoc gap.
\textit{(2) Outlier threshold:} instead of the paper's fixed threshold of 150, we use Otsu's method on the
pre-LayerNorm norm histogram, which generalises across model sizes.
\textit{(3) Position prediction:} we coarsen the $37{\times}37$ grid to $7{\times}7$ and subsample to
50k/20k patches to make logistic regression tractable.

\subsubsection{Experimental setup and code}
All experiments use PyTorch forward hooks to extract intermediate activations without modifying the model.

\textbf{Object discovery (Table~\ref{tab:object_discovery}).} We evaluate whether registers improve unsupervised
object localisation by running LOST on six model configurations across three datasets. A hook on
\texttt{model.blocks[-1].attn.qkv} (timm) or \texttt{resblocks[-1].attn} (OpenCLIP) captures the QKV
output; keys or values are separated by reshaping the $(B,N,3C)$ tensor and stripping special tokens by
index. Patch features pass through: L2 normalise $\to$ $G{=}\hat{K}\hat{K}^\top{+}b$, subtract per-row
mean $\to$ binarise at $\tau{=}0$ $\to$ degree computation $\to$ $\arg\min$ seed $\to$ seed expansion.
The tightest box is scored against GT; correct if $\max\,\text{IoU}{\geq}0.5$. Inference in FP16 on GPU
(batch~32); LOST per-image on CPU in float32. \textit{Metric: CorLoc:}— \% images with IoU(predicted box, best GT box) $\geq$ 0.5.\\
\textbf{Qualitative attention maps (Figure~\ref{fig:qualitative_attention}).} We visualise what CLS and register tokens
attend to, to verify that registers develop distinct spatial attention patterns. A hook on
\texttt{model.blocks[-1].attn} recomputes $\text{softmax}(QK^\top/\sqrt{d})$; the $(B,H,N,N)$ matrix
is averaged across heads and rows for CLS and each register are reshaped to $37{\times}37$ grids.
\textit{Metric: qualitative visual assessment.}\\
\textbf{LOST intermediate maps (Figures~\ref{fig:lost_intermediate_steps}, ~\ref{fig:seed_expansion_openclip}).} We inspect the three
intermediate LOST outputs on a single image to understand how registers affect each stage of the algorithm.
Same hook and pipeline as above; three maps are retained: inverse-degree, seed similarity row, and binary
expansion mask. For OpenCLIP values, the hook moves to the full \texttt{ResidualAttentionBlock} output —
pre-LN normalisation erases norm differences before the value projection.
\textit{Metric: qualitative assessment of intermediate maps.}\\
\textbf{Norm distributions (Figure~\ref{fig:token_norms_dinov2}).} We measure pre-LayerNorm token norms to verify
that registers absorb the high-norm outlier patches present in models without registers. A hook on
\texttt{model.norm} captures \texttt{inputs[0]} (pre-LN), not \texttt{output} — post-LN activations
destroy the bimodal signal. Token norms are split by type (CLS at 0, registers at 1--4, patches at 5--end)
and plotted as strip plots over 500 COCO images. \textit{Metric: visual comparison of norm distributions.}\\
\textbf{Information held by tokens (Tables~\ref{tab:outlier_token_probes}, ~\ref{tab:local_information_probes}).} We probe individual tokens
to quantify how much global and local information outlier patches and registers hold relative to normal
patches. A hook on \texttt{model.encoder.layer[-1]} (HuggingFace) captures post-LN features; the same
block's pre-LN input is used for outlier detection. Three probes per token: logistic regression on
FGVC-Aircraft (global), logistic regression on a $7{\times}7$ position grid (local), and Ridge regression
to raw pixels (reconstruction). \textit{Metric: top-1 accuracy; top-1 position prediction; mean L2 error.}\\
\textbf{Positional focus (Figure~\ref{fig:positional_focus}).} We visualise average attention patterns of CLS,
registers, and a centre patch to verify that registers adopt a global focus while patch tokens remain local.
\texttt{fused\_attn=False} forces softmax matrix materialisation; a hook on
\texttt{model.blocks[-1].attn.attn\_drop} captures the $(B,H,N,N)$ tensor. Rows for CLS (0),
registers (1--4), and centre patch (5+685) are averaged across $H{=}16$ heads over 500 Oxford-IIIT Pet
images and reshaped to $37{\times}37$. \textit{Metric: visual comparison of attention heatmaps.}

\subsubsection{Computational requirements}
\textbf{Hardware:} Single NVIDIA GPU (31.74 GiB VRAM) and single/double T4 (15.5 GiB each) on Kaggle. CPU work on MacBook Pro M3 Pro.
Total GPU time per experiment is reported in Table~\ref{tab:computationaltime}.
